\chapter{理論量の測定対応}
\label{ch:instantiation}

\section{本章の目的}

第\ref{part:theory}部で導入した記号のうち、どれが実際に測定できたのか、
測定できた場合にどの程度の値をとるのかを一覧する。
第\ref{part:results}部の各章はそれぞれ別個の領域を扱うため、
理論量との対応が散逸しやすい。本章はその索引を与える。

\textbf{測定できなかった量を明示することが本章の主目的である。}
枠組みが定義した量のうち、実データに接続できたものは半数に満たない。

\section{基本集合の実体}

\begin{center}
\begin{tabularx}{\textwidth}{@{}llX@{}}
\toprule
記号 & 領域 & 実現した集合と濃度 \\
\midrule
$N$ & 族2 & $\{$発行者, 購入者, 加盟店$\}$。発行者は 2,034 者（第三者型829・自家型1,205） \\
$N$ & 族5 & $\{$紹介事業者, 求職者, 求人企業$\}$。事業所 30,561、就職実績あり 13,946 \\
$N$ & 業界判定 & 法人企業統計の母集団。62業種 $\times$ 16資本金階層 \\
$T$ & 族2 & $\{$R2, R3, R4, R5, R6$\}$、$|T|=5$ \\
$T$ & 族5 & 令和6年度単年、パネルは令和2--7年度で $|T|=6$ \\
$T$ & 業界判定 & 2000--2024年度、$|T|=25$ \\
$\mathcal{D}$ & 族2 & 前払式支払手段の媒体4種（紙・磁気・IC・サーバ） \\
$\Omega$ & 全領域 & \textbf{観測されていない}。第\ref{part:unmanned}章で到来率のみ代理測定 \\
$\mathcal{F}_t$ & 全領域 & \textbf{観測されていない} \\
\bottomrule
\end{tabularx}
\captionof{table}{基本集合 $N$、$T$、$\mathcal{D}$、$\Omega$、$\mathcal{F}_t$ の各領域における実現状況。}
\label{tab:basic-sets}
\end{center}

\begin{remark}[$\Omega$ が測れないことの帰結]
状態空間そのものは観測できない。したがって
$\pi$ や $\delta$ の $\omega$ 依存性を直接検証することはできず、
本稿の実証はすべて $\omega$ について周辺化した量を扱っている。
唯一の例外は第\ref{part:unmanned}章で、
未想定状態の到来率 $\lambda_{\rm novel}$ を脆弱性件数で代理した。
\end{remark}

\section{写像 \texorpdfstring{$\Phi$}{Phi} の観測範囲}

第\ref{part:catalog}部で28類型を列挙したが、実データに接続できたのは以下に限られる（表\ref{tab:phi-observed-types}）。

\begin{center}
\begin{tabularx}{\textwidth}{@{}llX@{}}
\toprule
類型 & 実体 & 観測できた量 \\
\midrule
2-4 プリペイド & 前払式支払手段 & $D$, $P$, $\kappa$, $\dbar-\delta$ の上限 \\
5-1 成果報酬 & 有料職業紹介 & $\pi/$件、$y$、$\mathcal{G}$ の範囲 \\
6-1 取引額手数料 & 決済・配信基盤 & $\phi_{\rm barg}$ の価格 \\
6-4 機会課金 & マーケットプレイス集客 & $T_{\rm acq}$ の価格 \\
1-3 後払単発 & 法人企業統計の掛売 & $\kappa_C>0$ \\
\bottomrule
\end{tabularx}
\captionof{table}{28類型のうち実データに接続できた5類型と観測できた量。}
\label{tab:phi-observed-types}
\end{center}

すなわち\textbf{28類型のうち5類型}である。
族3（資産の時間貸し）、族4（補完財回収）、族7（支払者分離）は
一切の実測を行っていない。

\section{累積と信用ポジション}

\subsection{族2における \texorpdfstring{$D$}{D}, \texorpdfstring{$P$}{P}, \texorpdfstring{$\kappa$}{kappa}}

第\ref{sec:phi}章の定義に令和6年度の値を代入する。単位は百万円。

\begin{align}
P(\text{R6}) &= 28{,}150{,}350 \qquad \text{（発行額。顧客からの決済）} \\
D(\text{R6}) &= 28{,}241{,}983 \qquad \text{（回収額。企業からの履行）} \\
-\kappa_C(\text{R6}) &= 2{,}895{,}327 \qquad \text{（未使用残高。顧客が企業に与信）}
\end{align}

式\eqref{eq:kappa}に照らすと $\kappa_C = -2.90$ 兆円である。
負であるから顧客が発行者に与信している。
式\eqref{eq:ccc}を適用すると

\[
\CCC = \frac{W}{r} = \frac{2{,}895{,}327}{28{,}150{,}350}\times 365 = 37.5\ \text{日}.
\]

\subsection{媒体別の \texorpdfstring{$\CCC$}{CCC}}

同じ計算を媒体別に行うと、二桁の差が現れる。

\begin{center}
\begin{tabular}{@{}lrrr@{}}
\toprule
媒体 & $-\kappa_C$（百万円） & $r$（百万円） & $\CCC$（日） \\
\midrule
磁気型 & 562,980 & 107,925 & 1,904 \\
紙型 & 1,127,149 & 486,732 & 845 \\
サーバ型 & 670,552 & 14,331,607 & 17 \\
IC型 & 534,645 & 13,224,087 & 15 \\
\bottomrule
\end{tabular}
\captionof{table}{媒体別の $-\kappa_C$、$r$、$\CCC$。}
\label{tab:media-ccc}
\end{center}

\subsection{法人企業統計における \texorpdfstring{$\kappa$}{kappa}}

式\eqref{eq:kappa}を企業部門全体に適用する。2024年度、全産業（除く金融保険業）。

\begin{align}
\kappa_C &= \mathrm{DSO} = 71\ \text{日相当（10億円以上）},\quad 35\ \text{日相当（1千万円未満）}\\
\kappa_S &= -\mathrm{DPO} = -44\ \text{日相当（10億円以上）},\quad -21\ \text{日相当（1千万円未満）}\\
\sum_i \kappa_i &= 27\ \text{日相当（10億円以上）},\quad 14\ \text{日相当（1千万円未満）}
\end{align}

\textbf{$\kappa$ の符号は規模と業種の両方に依存し、業種の効果が支配的である。}
小売業では10億円以上で $\sum_i\kappa_i = -16$ 日と符号が反転する。

\section{成長制約の実測}
\label{sec:gstar-measured}

系\ref{cor:gstar}の $g^\star=(m+d-I/r)/\CCC$ は四つの量を要する。
本節は二段階に分けて測定する。

\begin{center}
\begin{tabularx}{\textwidth}{@{}lXl@{}}
\toprule
対象 & 用いる量 & 期間 \\
\midrule
長期の趨勢 & $m$ と $\CCC$ のみ。比 $m/\CCC$ を算出 & 2000--2024年度 \\
水準の測定 & $m$、$d$、$I$、$\CCC$。$g^\star$ を算出 & 2016--2024年度 \\
\bottomrule
\end{tabularx}
\captionof{table}{成長制約の実測における二段階の対象と期間。}
\label{tab:gstar-measurement-stages}
\end{center}

長期側で $d$ と $I$ を用いないのは、
減価償却費と固定資産残高を1990年代まで遡って取得していないためである。
$m$ は営業利益率を用いる。
第\ref{sec:growth}章の $m$ は運転資本の増加を賄う原資であるため、
販管費控除後の値が定義に近い。

\subsection{企業部門全体の推移（$m/\CCC$）}

\begin{center}
\begin{tabular}{@{}rrrr@{}}
\toprule
年度 & $\CCC$（日） & $m$ & $m/\CCC$ \\
\midrule
2000 & 24.9 & 2.62\% & 38.5\% \\
2005 & 24.4 & 3.16\% & 47.3\% \\
2009 & 28.0 & 2.01\% & 26.2\% \\
2015 & 30.6 & 3.95\% & 47.1\% \\
2020 & 35.4 & 3.06\% & 31.5\% \\
2024 & 37.2 & 5.01\% & 49.2\% \\
\bottomrule
\end{tabular}
\captionof{table}{企業部門全体の $\CCC$、$m$、$m/\CCC$ の推移（全産業）。}
\label{tab:mccc-trend}
\end{center}
\begin{center}\footnotesize 全産業（除く金融保険業）、全規模。$d$ と $I$ を含まない\end{center}

\begin{figure}[htbp]
\centering
\insertfig{gstar-trend}
\caption{$m/\CCC$ の推移（全産業、営業利益率ベース）。系\ref{cor:gstar}の $g^\star$ とは $d$ と $I$ の分だけ異なる。}
\label{fig:gstar-trend}
\end{figure}

\begin{proposition}[$m/\CCC$ は低下していない]
\label{prop:gstar-flat}
2000年度から2024年度にかけて、
$\CCC$ は 24.9 日から 37.2 日へ 49\% 悪化した。
しかし $m$ が 2.62\% から 5.01\% へ 91\% 改善したため、
\[
\frac{m}{\CCC} : \quad 38.5\% \;\longrightarrow\; 49.2\%
\]
となり、比はむしろ上昇した。
25年間の平均は 41.1\%、標準偏差は 7.3 ポイントである。
前半12年の平均 38.0\% に対し後半13年は 44.0\% であり、
低下傾向は認められない。
\end{proposition}

\begin{remark}[本命題は $g^\star$ そのものではない]
\label{rem:gstar-flat-caveat}
系\ref{cor:gstar}の $g^\star$ は $(m+d-I/r)/\CCC$ である。
命題\ref{prop:gstar-flat}が扱うのは $d$ と $I$ を含まない $m/\CCC$ にすぎない。

長期の時系列で $m/\CCC$ を用いるのは、
減価償却費と固定資産残高を1990年代まで遡って取得していないためである。
$d$ と $I$ を含む $g^\star$ は第\ref{sec:gstar-measured}節で
2016--2024年度について算出する。

したがって本命題は\textbf{比の趨勢についての主張}であり、
自己金融可能成長率の水準についてではない。
$I$ が趨勢的に変化していれば、$g^\star$ の推移は $m/\CCC$ と異なりうる。
\end{remark}

\begin{remark}[分子と分母を分けて測る必要]
$\CCC$ の悪化のみを見れば成長制約の悪化が示唆される。
しかし系\ref{cor:gstar}は比であり、
\textbf{分母の変化のみから結論を導くことはできない}。
$m$ を取得せずに $\CCC$ のみで論じると、事実と逆の結論に至る。

$m/\CCC$ の変動を駆動しているのは主に $m$ である。
谷は2008--2009年度（26\%台）と2020年度（31.5\%）にあり、
いずれも $\CCC$ ではなく利益率の低下による。
\end{remark}

\subsection{減価償却と投資を含めた \texorpdfstring{$g^\star$}{g*}}

系\ref{cor:gstar}は $m$ に加えて減価償却費率 $d$ と投資 $I$ を要する。
注\ref{rem:no-steady-investment}のとおり $I\approx d\,r$ の仮定は置かない。
$I$ は固定資産残高の差分に減価償却費を加えて推定した。
土地は減価償却の対象外であるため除いている。
\begin{equation}
I \approx \bigl(\text{有形}+\text{無形}-\text{土地}\bigr)_t
 - \bigl(\text{有形}+\text{無形}-\text{土地}\bigr)_{t-1} + \text{減価償却費}_t
\label{eq:I-estimate}
\end{equation}

\begin{center}
\begin{tabular}{@{}lrrrrrrr@{}}
\toprule
規模（資本金） & $\CCC$ & $m$ & $d$ & $I/r$ & $g^\star$(2024) & 9年平均 & SD \\
\midrule
10億円以上 & 47 & 7.8\% & 3.1\% & 4.2\% & 52.0\% & 49.3\% & 6.0 \\
1億円以上 -- 10億円未満 & 35 & 4.8\% & 2.1\% & 4.3\% & 27.1\% & 39.6\% & 12.2 \\
5千万円以上 -- 1億円未満 & 30 & 3.9\% & 1.8\% & 1.9\% & 45.8\% & 33.0\% & 15.2 \\
2千万円以上 -- 5千万円未満 & 25 & 3.5\% & 2.1\% & 1.2\% & 64.3\% & 35.0\% & 17.0 \\
1千万円以上 -- 2千万円未満 & 37 & 1.8\% & 2.3\% & 5.5\% & \textminus 13.8\% & 15.5\% & 30.8 \\
1千万円未満 & 26 & 1.5\% & 2.9\% & 2.0\% & 34.7\% & \textbf{3.1\%} & 41.4 \\
\bottomrule
\end{tabular}
\captionof{table}{資本金階層別の $\CCC$、$m$、$d$、$I/r$、$g^\star$（2024年度、9年平均、SD）。}
\label{tab:gstar-by-size}
\end{center}
\begin{center}\footnotesize 全産業（除く金融保険業）、$\CCC$ は日、SD は $g^\star$ の9年標準偏差。
取得したデータは2015--2024年度の10年だが、式\eqref{eq:I-estimate}の $I$ が前年の固定資産残高との
差分を要するため $g^\star$ を算出できるのは2016年度以降であり、平均と SD は2016--2024年度の9年で取る\end{center}

\begin{remark}[単年値は規模により信頼性が異なる]
\label{rem:gstar-reliability}
$I/r$ の9年標準偏差は、10億円以上で 0.3、1千万円未満で 2.9 と一桁違う。
式\eqref{eq:I-estimate}は残高の差分を用いるため、
除却・売却と\textbf{資本金階層の移動}が混入する。
資本金は増資により業績と無関係に変わるため、階層間の移動が生じる。

移動の向きは一定でない。1千万円未満の $I/r$ は
2016年に $-1.6\%$、2022年に $+7.5\%$ を記録しており、
一方向のバイアスではない。したがって平均を取れば影響は概ね相殺されるが、
単年値は小規模階層では信頼できない。

本表が単年と9年平均を併記し標準偏差を付すのはこのためである。
\end{remark}

注\ref{rem:gstar-reliability}を踏まえ、9年平均で判定する。

\begin{proposition}[小規模層の制約は $\CCC$ ではない]
\label{prop:small-margin}
9年平均で見ると、1千万円未満の $g^\star$ は 3.1\%、
10億円以上の 49.3\% に対し一桁以上小さい。
この差は $\CCC$ に由来しない。
小規模層の $\CCC$（26--37日）はむしろ大企業（47日）より短い。

原因は分子である。$m+d-I/r$ は
10億円以上で 6.7\%、1千万円未満で 2.4\% であり、
$m$ が 1.5\% と低いことに加え、$d$ を上回る投資負担が寄与している。
\end{proposition}

単年では 1千万円未満が $34.7\%$ と高く出るが、
標準偏差 41.4 を考慮すればこの値に意味を読むべきではない。

\subsection{業種別の \texorpdfstring{$g^\star$}{g*}}

\begin{center}
\begin{tabular}{@{}lrrrrrr@{}}
\toprule
業種 & $\CCC$ & $m$ & $d$ & $I/r$ & 9年平均 $g^\star$ & SD \\
\midrule
製造業 & 42 & 5.4\% & 2.8\% & 4.2\% & 34.8\% & 5.6 \\
全産業 & 37 & 5.0\% & 2.5\% & 3.5\% & 37.7\% & 7.5 \\
卸売業 & 33 & 2.2\% & 0.7\% & 0.4\% & 20.9\% & 6.9 \\
建設業 & 40 & 4.5\% & 1.5\% & 1.4\% & 47.1\% & 11.3 \\
小売業 & 23 & 2.7\% & 1.5\% & 1.6\% & 29.0\% & 12.2 \\
職業紹介・労働者派遣業 & 33 & 3.9\% & 0.6\% & \textminus 0.3\% & 42.7\% & 13.9 \\
情報通信業 & 58 & 8.2\% & 4.8\% & 5.2\% & 61.2\% & 16.5 \\
その他のサービス業 & 34 & 5.4\% & 3.0\% & 5.5\% & 47.9\% & 32.3 \\
\midrule
宿泊業 & 22 & 6.6\% & 3.5\% & 6.3\% & \textminus 204.2\% & \textbf{872} \\
飲食サービス業 & 5 & 1.2\% & 2.4\% & \textminus 0.4\% & 266.9\% & \textbf{3334} \\
\bottomrule
\end{tabular}
\captionof{table}{業種別の $\CCC$、$m$、$d$、$I/r$、9年平均 $g^\star$（SD付き）。}
\label{tab:gstar-by-industry}
\end{center}

表\ref{tab:gstar-by-industry}の下二業種は注\ref{rem:gstar-divergence}の発散が生じている。
飲食サービス業の $\CCC$ は 5 日であり、分子が 1 ポイント動けば
$g^\star$ は約 73 ポイント動く。標準偏差 3334 はこの増幅の帰結であり、
\textbf{$g^\star$ が指標として機能していない}。

\begin{remark}[$\CCC$ が短いことの二つの意味]
第\ref{part:catalog}部の類型でいえば、飲食サービス業は族1-1（即時交換）であり、
$\tau_i\approx 0$ が構造的な性質である（系\ref{cor:ccc-level}）。
$\CCC$ が短いこと自体は資金繰りの健全性を示すが、
\textbf{$g^\star$ という比の指標は同時に無意味になる}。
系\ref{cor:gstar}が有用なのは $\CCC$ が数十日以上の対象に限られる。
\end{remark}

\subsection{小規模層の \texorpdfstring{$m$}{m} は役員報酬で決まる}
\label{sec:officer-comp}

命題\ref{prop:small-margin}は $m$ が制約であることを示したが、
$m$ が低い理由までは述べていない。費用構造を分解する。

\begin{center}
\begin{tabular}{@{}lrrrr@{}}
\toprule
規模（資本金） & 粗利率 & 販管費率 & 人件費率 & $m$ \\
\midrule
10億円以上 & 23.6\% & 15.8\% & 7.7\% & 7.8\% \\
1億円以上 -- 10億円未満 & 21.5\% & 16.8\% & 10.3\% & 4.8\% \\
5千万円以上 -- 1億円未満 & 25.9\% & 22.0\% & 12.3\% & 3.9\% \\
2千万円以上 -- 5千万円未満 & 26.0\% & 22.4\% & 15.0\% & 3.5\% \\
1千万円以上 -- 2千万円未満 & 33.8\% & 32.0\% & 17.7\% & 1.8\% \\
1千万円未満 & \textbf{46.3\%} & \textbf{44.8\%} & 22.7\% & 1.5\% \\
\bottomrule
\end{tabular}
\captionof{table}{資本金階層別の粗利率・販管費率・人件費率と $m$（2024年度）。}
\label{tab:cost-structure-by-size}
\end{center}
\begin{center}\footnotesize 全産業（除く金融保険業）、2024年度\end{center}

\textbf{小規模層は粗利率が高い。} 大企業の 23.6\% に対し 46.3\% である。
一方で販管費率は 15.8\% に対し 44.8\% と大きい。
注\ref{rem:cadmin-negligible}のとおり、
この差の一部は $C_{\mathrm{admin}}$ である可能性がある。
請求と督促が人手で行われる規模では、契約の運用費用が販管費に計上される。
本稿はこれを分離していない。
すなわち $m$ の低さは価格支配力の欠如によるものではない。
落ちているのは販管費の段階である。

人件費を売上原価側に戻すと、粗利率の差は 22.7 ポイントから 7.7 ポイントへ縮む
（15.9\% 対 23.6\%）。差の三分の二は会計上の分類差であった。
ただし調整後もなお小規模層のほうが高い。

さらに人件費を役員と従業員に分解すると、原因が特定できる。

\begin{center}
\begin{tabular}{@{}lrrrr@{}}
\toprule
規模（資本金） & 役員報酬率 & 従業員給与率 & 役員比率 & $m$ \\
\midrule
10億円以上 & 0.16\% & 7.6\% & 2.0\% & 7.8\% \\
1億円以上 -- 10億円未満 & 0.40\% & 9.9\% & 3.9\% & 4.8\% \\
5千万円以上 -- 1億円未満 & 0.99\% & 11.3\% & 8.1\% & 3.9\% \\
2千万円以上 -- 5千万円未満 & 2.35\% & 12.7\% & 15.6\% & 3.5\% \\
1千万円以上 -- 2千万円未満 & 3.98\% & 13.7\% & 22.5\% & 1.8\% \\
1千万円未満 & \textbf{8.04\%} & 14.6\% & \textbf{35.5\%} & 1.5\% \\
\bottomrule
\end{tabular}
\captionof{table}{資本金階層別の役員報酬率・従業員給与率・役員比率と $m$。}
\label{tab:officer-comp-by-size}
\end{center}
\begin{center}\footnotesize 役員比率は人件費に占める役員報酬の割合\end{center}

\begin{proposition}[役員報酬が $m$ の規模差を説明する]
\label{prop:officer}
売上高に対する役員報酬の比率は、資本金10億円以上で 0.16\%、
1千万円未満で 8.04\% であり、\textbf{50倍の差}がある。
最大規模層と最小規模層の $m$ の差は 6.3 ポイントであるのに対し、
役員報酬率の差は 7.88 ポイントであり、前者を説明しうる大きさである。

役員報酬を営業利益に戻すと単調な序列は消失する。
\begin{center}
\begin{tabular}{@{}lrr@{}}
\toprule
規模 & $m$ & 役員報酬を戻した後 \\
\midrule
10億円以上 & 7.8\% & 8.0\% \\
2千万円以上 -- 5千万円未満 & 3.5\% & 5.9\% \\
1千万円未満 & 1.5\% & \textbf{9.5\%} \\
\bottomrule
\end{tabular}
\captionof{table}{役員報酬を営業利益に戻した前後の $m$（規模別）。}
\label{tab:officer-comp-adjusted}
\end{center}
最小規模層が最も高くなる。
業種内比較でも同様であり、$m$ が小規模層で高い業種は9業種中1つであるのに対し、
役員報酬を戻すと7業種に逆転する。
\end{proposition}

\begin{remark}[$m$ の定義との齟齬]
第\ref{sec:growth}章では $m$ を「運転資本の増加を賄う原資」と定義した。
役員報酬が裁量的であるならば、小規模事業者が実際に使える原資は
営業利益ではなく\textbf{営業利益と役員報酬の和}である。
その基準では 1千万円未満層の分子は $m+d-I/r$ の 2.4\% ではなく、
役員報酬 8.04\% を加えた 10.4\% となる。
$\CCC=26$ 日で除せば 146\% であり、
大企業の 51.7\%（第\ref{sec:gstar-measured}節、9年平均 49.3\%）を大きく上回る。

すなわち命題\ref{prop:small-margin}の「小規模層は $m$ に拘束される」は、
公表される営業利益を $m$ とみなす限りで正しいが、
\textbf{定義に忠実な $m$ を用いれば結論が反転する}。
\end{remark}

\subsection{個人企業との比較}
\label{sec:individual-firms}

命題\ref{prop:officer}は法人企業統計に基づく。
異なる統計、異なる規模区分、かつ\textbf{法人格を持たない主体}を含む標本で
同じ構造が現れるかを確認する。

中小企業実態基本調査を用いた。標本約11万社、
従業者規模により法人4区分と個人企業が同一表に並ぶ。
個人企業には役員報酬という費目が存在せず、
事業主の労働対価は営業利益に残る。

\begin{quote}
\textbf{予測U}：個人企業の営業利益率は、
法人小規模層の「役員報酬を戻した後」の値に近い。
両者は実質的に同じ経済主体であり、会計上の区分が異なるにすぎない。
\end{quote}

\begin{center}
\begin{tabular}{@{}lrrrrrrrr@{}}
\toprule
区分 & 2019 & 2020 & 2021 & 2022 & 2023 & 2024 & 平均 & SD \\
\midrule
法人 5人以下 & 2.7 & 1.1 & 2.6 & 3.0 & 2.6 & 3.0 & 2.5 & 0.7 \\
法人 6--20人 & 2.4 & 1.5 & 2.5 & 2.8 & 3.0 & 3.2 & 2.6 & 0.6 \\
法人 21--50人 & 2.8 & 2.1 & 3.2 & 2.8 & 3.4 & 3.7 & 3.0 & 0.6 \\
法人 51人以上 & 3.4 & 2.8 & 3.2 & 3.6 & 3.6 & 4.0 & 3.4 & 0.4 \\
個人企業 & 17.0 & 18.5 & 18.2 & 15.8 & 15.4 & 16.3 & \textbf{16.9} & 1.3 \\
\bottomrule
\end{tabular}
\captionof{table}{法人規模別・個人企業の営業利益率の推移（2019--2024年度）。}
\label{tab:individual-vs-corp-margin}
\end{center}
\begin{center}\footnotesize 営業利益率（\%）、全産業、決算年度\end{center}

\textbf{予測Uは支持される。} 個人企業の営業利益率は6年間すべてで法人の5倍前後であり、
標準偏差1.3と安定している。単年の特殊事情ではない。
第\ref{sec:officer-comp}節で法人の役員報酬を戻した値は 9.5\% であった。
個人企業の 16.9\% はこれを上回るが、
戻す前の 1.5--3.0\% よりはるかに近い。

人件費率がこの差を説明する。

\begin{center}
\begin{tabular}{@{}lrrr@{}}
\toprule
区分 & 営業利益率 & 人件費率 & 合計 \\
\midrule
法人 5人以下 & 2.5\% & 14.7\% & 17.2\% \\
法人 6--20人 & 2.6\% & 16.9\% & 19.5\% \\
個人企業 & 16.9\% & \textbf{9.8\%} & 26.7\% \\
\bottomrule
\end{tabular}
\captionof{table}{法人・個人企業の営業利益率と人件費率（6年平均）。}
\label{tab:individual-vs-corp-laborcost}
\end{center}
\begin{center}\footnotesize 6年平均。人件費率は売上原価の労務費と販管費の人件費の和\end{center}

\textbf{個人企業の人件費率が最も低い。}
事業主の労働が費用に計上されないためである。
法人では役員報酬として人件費に入る分が、個人企業では営業利益に残る。

\begin{remark}[コロナ期に符号が逆転する]
\label{rem:covid-reversal}
2020年度、法人は全階層で営業利益率が低下した（5人以下で 2.7\% $\to$ 1.1\%）。
\textbf{個人企業のみ上昇している}（17.0\% $\to$ 18.5\%）。

これは注\ref{rem:officer-undecidable}で述べた観測と同型である。
個人企業の営業利益は事業主の取り分を含むため、
売上が減少しても事業主が取り分を減らさなければ比率は低下しない。
従業員給与や外注費が先に削られれば、むしろ上昇する。

法人の役員報酬が利益変動と独立に動くことと、
個人企業の営業利益率が不況期に上昇することは、
\textbf{同一の現象の二つの現れ}である。
\end{remark}

\begin{remark}[個人企業の設備投資]
\label{rem:individual-capex}
同調査により個人企業の $I/r$ が得られる。6年平均 1.9\%、標準偏差 1.0。
法人 51人以上の 3.1\%（SD 0.1）に対し低く、かつ不安定である。
設備投資を行った企業の割合は 8.1\%（2024年度）であり、9割以上が投資を行っていない。

第\ref{ch:solo}章では一人事業について $A\approx 0$ を仮定した。
この仮定は個人企業の実データと整合する。
\end{remark}

\begin{remark}[判定できないこと]
\label{rem:officer-undecidable}
役員報酬が裁量的な利益処分なのか労働の対価なのかは、本稿では判定できない。

裁量性を支持する観測として、1千万円未満層の役員報酬率は
2015--2024年度を通じて 7.6--9.1\% の狭い範囲にとどまる一方、
同期間の $m$ は $-2.0\%$ から $+2.1\%$ まで変動している。
\textbf{利益が変動しても役員報酬は動かない。}

しかし労働対価説も棄却できない。小規模企業の役員は実質的に労働者として稼働しており、
報酬がその対価である可能性は同程度にある。
両者の区別には役員一人あたり報酬額と労働時間が必要だが、
\textbf{役員は労働基準法上の労働者ではないため労働時間の統計が存在しない}。
毎月勤労統計調査も賃金構造基本統計調査も対象を労働者に限っている。
これは第\ref{sec:obstacle-taxonomy}節の（i）測定の不在に属する。
\end{remark}

\subsection{業種別（2024年度）}

\begin{center}
\begin{tabular}{@{}lrrr@{}}
\toprule
業種 & $\CCC$（日） & $m$ & $g^\star$ \\
\midrule
純粋持株会社 & 12 & 57.7\% & 1710\% \\
水運業 & 7 & 4.8\% & 261\% \\
鉱業、採石業、砂利採取業 & 29 & 19.0\% & 236\% \\
自動車・同附属品製造業 & 9 & 5.1\% & 201\% \\
\midrule
リース業 & 210 & 6.6\% & 11.4\% \\
石油製品・石炭製品製造業 & 20 & 0.6\% & 11.4\% \\
繊維工業 & 78 & 2.0\% & 9.2\% \\
漁業 & 40 & \textminus 1.5\% & 負 \\
農業、林業 & 22 & \textminus 4.0\% & 負 \\
\bottomrule
\end{tabular}
\captionof{table}{業種別の $\CCC$、$m$、$g^\star$（2024年度、上位・下位の例）。}
\label{tab:gstar-industry-extremes}
\end{center}

分子 $m+d-I/r$ が負の業種では $g^\star<0$ となり、
命題\ref{prop:fcf}より\textbf{成長率がゼロでもフリーキャッシュフローが負}になる。

$\CCC$ が長いことと $g^\star$ が小さいことは一致しない（表\ref{tab:gstar-industry-extremes}）。
リース業は $\CCC=210$ 日と最長だが $m=6.6\%$ が高いため、
$m$ のみで比較すれば石油製品・石炭製品製造業（$\CCC=20$ 日、$m=0.6\%$）と同水準になる。
\textbf{$\CCC$ 単独では成長制約を判定できない。}
なお $d$ と $I$ を含めた $g^\star$ は第\ref{sec:gstar-measured}節に示す。

\section{権利と行使、認知余剰}

式\eqref{eq:cog}の $\phi_{\rm cog}=p\cdot\E[\dbar-\delta]$ に対応する測定値は、
族2の有効期限到来等による回収額のみである。

\begin{align}
p\cdot\E[\dbar-\delta] &\leq 20{,}634\ \text{百万円}\\
\frac{20{,}634}{2{,}895{,}327} &= 0.71\%\quad\text{（未使用残高比）}\\
\frac{20{,}634}{28{,}150{,}350} &= 0.073\%\quad\text{（発行額比）}
\end{align}

不等号であるのは、この金額に払戻しによる回収が混入しているためである。
すなわち\textbf{本稿の測定では $\phi_{\rm cog}$ の上界のみが得られており、値は未知である}。

\subsection{\texorpdfstring{$\phi_{\rm cog}$}{phi-cog} は測定可能である}
\label{sec:phicog-measurable}

先行研究には直接の測定例がある。

\cite{dellavigna2006} は米国のフィットネスクラブ三施設、
会員 7,752 名の契約選択と日次来館記録を突き合わせ、
式\eqref{eq:cog}の $\dbar-\delta$ を直接測定している。
会員の予測来館回数は月 9.50 回、実績は 4.17 回であった。
月額 \$70 超の会員は平均 4.3 回しか来場せず、1回あたり \$17 以上を支払っている。
10回券を用いれば \$10 で済む。会員期間を通じた損失は平均 \$600、
総支払額は約 \$1{,}400 であるから、$\phi_{\rm cog}$ は支払額の約 43\% にあたる。

\cite{einav2025} はより広い範囲を扱う。
決済カードネットワークの網羅データを用い、
カード更新月には能動的な更新判断が強制されることを外生ショックとして利用した。
10のサブスクリプションについて、\textbf{不注意が売り手の収益を平均 89\% 押し上げる}
（推定幅は 14\% から 200\% 超）。
注意パラメータ $\lambda$ の推定値は平均 0.18、サービス間で 0.04 から 0.50 である。

さらに \cite{nber2026pnbl} は族2-4（プリペイド）を直接扱う。
飲食店で使える前払い残高にボーナスを付す仕組みについて、
400万人超の取引データを分析し、
\textbf{前払い価値の約40\%が使用されない}ことを記録している。
事業者はボーナス1ドルあたり中央値5.5ドルの退蔵利益を得ており、
支出増による収益よりも退蔵からの利益のほうが大きい。

\begin{proposition}[$\phi_{\rm cog}$ の測定に要する三条件]
\label{prop:phicog-requirements}
$\phi_{\rm cog}$ を推定するには次の三つが同時に必要である。
\begin{enumerate}[label=(\arabic*)]
\item \textbf{契約上の権利 $\dbar$}：料金体系と利用上限
\item \textbf{実現利用 $\delta$}：同一主体の個票記録。集計値では不可
\item \textbf{注意を強制する外生イベント}：カード更新、規約変更、価格改定など
\end{enumerate}
\end{proposition}

\begin{proof}[根拠]
(1)(2) がなければ差 $\dbar-\delta$ が定義できない。
(3) がなければ、観測された継続が「価値があるから継続した」のか
「注意を払わなかったから継続した」のかを区別できず、
$\lambda$ が識別されない。
\cite{dellavigna2006} は (1)(2) を契約データと入館記録で、
\cite{einav2025} は (3) をカード更新で満たしている。
\end{proof}

本稿の族2の測定で $\phi_{\rm cog}$ が得られなかったのは、
\textbf{(2) と (3) を欠いていたためである}。
集計統計には個票がなく、外生ショックも存在しない。
第\ref{sec:obstacle-taxonomy}節の分類でいえば、
これは（i）測定の不在ではなく（iv）アクセス制約である。
測定手法は確立しており、必要なのはデータへの到達である。

\subsection{メニューからの迂回：予測Kの検証}
\label{sec:phicog-detour}

命題\ref{prop:phicog-requirements}の条件(2)個票の実現利用は、
公表統計では得られない。しかし\textbf{定義を置き換えることで迂回できる場合がある}。

定額 $c$ と都度 $p$ が同一メニューに併存するとき、
損益分岐となる利用回数が計算できる。
\begin{equation}
\delta^{*} = c / p
\label{eq:delta-star}
\end{equation}
$\dbar$ が「無制限」であって測定できない場合でも、
式\eqref{eq:delta-star}は\textbf{メニューのみから求まる}。
これを $\dbar$ の代理として用いれば、
必要なのは業種平均の利用回数 $\delta$ だけになる。

\begin{quote}
\textbf{予測K}：定額会員について $\delta^{*} - \delta > 0$ である。
すなわち平均的な会員は損益分岐に達していない。
\end{quote}

\paragraph{結果。}
$\delta$ は特定サービス産業実態調査より、個人会員1人あたり年89回、
すなわち月 7.42 回である。
$\delta^{*}$ は総合型大手一社の実料金（2026年、税込）から算出した。

\begin{center}
\begin{tabular}{@{}lrrrr@{}}
\toprule
施設カテゴリ & 無制限プラン & 都度利用 & $\delta^{*}$ & $\delta^{*}-\delta$ \\
\midrule
I & 12,680 & 2,530 & 5.01 & \textminus 2.40 \\
II & 13,980 & 2,860 & 4.89 & \textminus 2.53 \\
III & 16,580 & 3,190 & 5.20 & \textminus 2.22 \\
IV & 19,080 & 3,630 & 5.26 & \textminus 2.16 \\
\bottomrule
\end{tabular}
\captionof{table}{施設カテゴリ別の無制限プラン・都度利用料金と損益分岐利用回数 $\delta^{*}$。}
\label{tab:predk-breakeven}
\end{center}

\textbf{予測Kは棄却される。} 四カテゴリすべてで符号が負であり（表\ref{tab:predk-breakeven}）、
平均的な会員にとって定額契約は合理的である。

会費総額を延べ利用者数で除した実効単価も同じ方向を示す。
2000年の1,147円から2023年の815円へ、25年間で29\%低下している。
都度利用の相場が2,530--3,630円であることに照らせば、
\textbf{会員は都度払いより安く利用している}。

\begin{remark}[迂回路の適用条件]
\label{rem:detour-limit}
式\eqref{eq:delta-star}が一意に定まるのは、メニューが定額と都度の二択である場合に限る。
実際の料金体系は階段状であり、上記の事業者は月4回、月8回、無制限の三段階を持つ。
このとき比較対象によって閾値が変わる。
無制限プランを選ぶ合理性は、都度利用を基準とすれば月5.0回、
月8回プランを基準とすれば月9.1--9.5回を要する。
\textbf{平均利用7.42回は、前者では合理的、後者では不合理となる。}

したがって本迂回路は二択メニューにのみ適用でき、
階段メニューでは $\dbar$ の代理が一意に定まらない。
\end{remark}

\begin{remark}[本節から $\phi_{\rm cog}=0$ は導けない]
予測Kの棄却は平均についての結果である。
\cite{dellavigna2006} の知見は分布に関するものであり、
平均が損益分岐を上回っていても下位層が損失を被る構造は排除されない。
平均と分布は別の主張である。
\end{remark}

\begin{remark}[本稿の枠組みで残る未測定領域]
先行研究の対象はサブスクリプションと会員制、
すなわち第\ref{part:catalog}部の族2-1 に集中している。
族2-4（プリペイド、ギフトカード）と族2-5（オプション・保証）について
同種の推定は見当たらない。
また $\phi$ 全体に占める $\phi_{\rm cog}$ の\emph{比率}を業種横断で推計した研究も
確認できていない。\cite{einav2025} が測ったのは収益の増分であり、
利益の構成比ではない。
\end{remark}

取引頻度 $\nu$ の代理変数としては、媒体別の購入事由が使える。

\begin{center}
\begin{tabular}{@{}lrrr@{}}
\toprule
媒体 & 自身で使用 & 贈答用 & $\CCC$（日） \\
\midrule
IC型 & 95.5\% & 0.0\% & 15 \\
サーバ型（オンライン） & 78.1\% & 14.0\% & 17 \\
紙型 & 41.5\% & 30.1\% & 845 \\
\bottomrule
\end{tabular}
\captionof{table}{媒体別の購入事由（自身で使用・贈答用）と $\CCC$。}
\label{tab:media-purchase-reason}
\end{center}

贈答比率が高いほど $\nu$ が低く、$\CCC$ が長い。
式\eqref{eq:decay}の $\lambda(\nu)$ そのものは測定していない。

\section{情報構造 \texorpdfstring{$\mathcal{G}$}{G} の実体}

第\ref{sec:info}章の $\mathcal{G}$（検証可能な共有情報）は、
族5において具体的な項目集合として観測できる。
人材サービス総合サイトに掲載される項目が $\mathcal{G}$ の実体である。

\begin{center}
\begin{tabularx}{\textwidth}{@{}lXl@{}}
\toprule
記号 & 実体 & 状態 \\
\midrule
$y$ & 無期雇用就職者のうち6か月以内離職者数 & 数値、6年分 \\
$\mathcal{G}$ の広さ & 「離職が判明せず」の人数（少ないほど $\mathcal{G}$ が広い） & 数値。観測2件とも全年度ゼロ \\
$b$（成果連動係数） & 返戻金制度 & 有無は二値。率と期間は形式が不統一 \\
価格側の代理 & 職種別手数料実績率 & 連続量。307社（第\ref{sec:predX-result}節） \\
$\sigma^2$（結果のノイズ） & 職種別のベースライン離職率 & \textbf{未測定} \\
\bottomrule
\end{tabularx}
\captionof{table}{情報構造 $\mathcal{G}$ の実体と各記号の観測状態（族5）。}
\label{tab:info-structure-observed}
\end{center}

$y$ の実測値は、確認できた2事業所で次のとおりである。

\[
\frac{y}{\text{無期就職者数}} = \frac{124}{4{,}348} = 2.85\%, \qquad
\frac{3}{189} = 1.59\% .
\]

例\ref{ex:bstar}の $b^\star=1/(1+r\sigma^2 k)$ を検証するには
$b$ と $\sigma^2$ の両方が連続量として必要だが、
返戻金の率と期間は形式が不統一であり、
$\sigma^2$ からは制御不能な部分を取り出せない。

ただし価格側からの検定は実施できた。
職種別手数料実績率は307社について連続量として得られ、
離職率との相関を測定した（第\ref{sec:predX-result}節）。

\section{余剰の三分解}

式\eqref{eq:surplus}の三成分を分離できるのは第\ref{sec:platform-fee}章の事例のみである。

\begin{center}
\begin{tabularx}{\textwidth}{@{}lrX@{}}
\toprule
成分 & 測定値 & 根拠 \\
\midrule
$\phi_{\rm prod}$（機能の対価） & 20.0 ポイント & Gumroad の同一プラットフォーム内比較。集客の有無のみが異なる \\
$\phi_{\rm barg}$（支配力レント） & 15.0 ポイント & App Store の規模別二段階。機能は同一 \\
$\phi_{\rm risk}$（移転と集約の対価） & --- & この二社では $\approx 0$。個別審査型では正 \\
$\phi_{\rm cog}$ & --- & 本領域では測定していない \\
\bottomrule
\end{tabularx}
\captionof{table}{余剰の三分解の測定値と根拠（プラットフォーム手数料の事例）。}
\label{tab:surplus-decomposition-measured}
\end{center}

これは\textbf{単一の取引について三分解のうち二成分を数値的に分離できた唯一の事例}である。
他の領域では成分の分離手続きを持たない。

対象の二社は公表料金表により一律の料率で運用する仲介者であり、
リスク成分 $\phi_{\rm risk}$ が無視できたために分離が可能であった。
料率が個別審査で決まる仲介者では第三項が加わり、分離は成立しない
（注\ref{rem:fee-incomplete}）。

\section{立ち上げ費用の各項}

式\eqref{eq:E-required-2}の各項について、外部化価格が測定された。
取引額 \$50 における実効料率の差分である。

\begin{center}
\begin{tabular}{@{}llr@{}}
\toprule
項 & 外部化の手段 & 価格 \\
\midrule
$T_{\rm inst}$ & 決済代行 & 3.6\% \\
$T_{\rm inst}$（税務） & Merchant of Record & +1.4〜2.4 ポイント \\
$T_{\rm dev}$（配信） & ホスティング付き販売基盤 & +8〜9 ポイント \\
$T_{\rm acq}$ & マーケットプレイス集客 & +20 ポイント \\
$T_{\rm credit}$ & 無料期間による自弁 & $\kappa>0$ のコスト \\
\bottomrule
\end{tabular}
\captionof{table}{立ち上げ費用の各項と外部化の手段・価格。}
\label{tab:launch-cost-externalization}
\end{center}

\textbf{$T_{\rm acq}$ が他の全項目の合計を上回る。}
$T_{\rm dev}$ そのもの（開発期間）と $\hat\Omega$ の列挙は外部化できないため価格を持たない。

\section{選択と識別に関する量}

\begin{center}
\begin{tabularx}{\textwidth}{@{}llX@{}}
\toprule
記号 & 測定値 & 出所 \\
\midrule
$\Prob(S=1)$（上場） & $3{,}975 / 3{,}600{,}000 = 0.11\%$ & 第\ref{sec:frame}節 \\
$\Prob(S=1)$（言及される） & $10^{-4}$ の桁 & 同上 \\
$\tau$（倒産時刻） & \textbf{未観測} & 生存条件づけにより標本外 \\
$(a,p,c)$ & \textbf{識別不能} & 命題\ref{prop:apc} \\
$\lambda_{\rm novel}$ の上界 & 48,185 件/年（CVE 公開） & 第\ref{part:unmanned}章 \\
$\lambda_{\rm novel}$ の実効値 & 245 件/年（悪用確認）$=0.51\%$ & 同上 \\
\bottomrule
\end{tabularx}
\captionof{table}{選択と識別に関する量の測定値と出所。}
\label{tab:selection-identification-measures}
\end{center}

\section{測定できなかった量の一覧}

以下は本稿の枠組みが定義したが、実データに接続できなかった量である（表\ref{tab:unmeasured-quantities}）。

\begin{center}
\begin{tabularx}{\textwidth}{@{}llX@{}}
\toprule
記号 & 名称 & 接続できない理由 \\
\midrule
$M(t)$ & 現金残高 & 個社データが必要。集計では意味を持たない \\
$\Omega,\ \mathcal{F}_t$ & 状態空間、フィルトレーション & 原理的に観測されない \\
$v$ & 評価写像 & 主観的。定義が固定できない \\
$\beta_i$ & 主体別の割引因子 & 実験的手法を要する \\
$\mathrm{Cap}$ & 容量上限 & 事業者ごとに異なり、開示されない \\
$\sigma^2$ & 結果のノイズ & 職種別の分散が未集計（第\ref{sec:what-to-measure}節） \\
$\lambda(\nu)$ & 認知余剰の減衰率 & 分子である $\phi_{\rm cog}$ 自体が未測定 \\
$\phi_{\rm cog}$ & 認知余剰 & 本稿では上界のみ。測定手法は確立（第\ref{sec:phicog-measurable}節） \\
$\hat\Omega$ & 想定状態空間 & 設計者の内部にのみ存在 \\
$\tau$ & 倒産時刻 & 生存条件づけにより標本外（命題\ref{prop:pathwise}） \\
& & \\
\bottomrule
\end{tabularx}
\captionof{table}{本稿の枠組みが定義したが実データに接続できなかった量の一覧。}
\label{tab:unmeasured-quantities}
\end{center}

$m$ は当初この一覧に含まれていたが、取得して第\ref{sec:gstar-measured}節で用いた。
$\phi_{\rm cog}$ も原理的には測定可能であり（命題\ref{prop:phicog-requirements}）、
本稿が到達できなかったのはデータの粒度による。
残りは（i）測定の不在または原理的な観測不能に属する。

\begin{remark}[枠組みの過剰さ]
定義した量の約半数が実データに接続できていない。
これは枠組みが世界より細かく記述しているということであり、
記述の精密さと検証可能性は独立であることを示す。
第\ref{sec:obstacle-taxonomy}節の四分類で言えば、
接続できない量の多くは（i）測定の不在に属し、
労力の投入では解決しない。
\end{remark}
